\subsection{半单李群}

命题 26. 设 \(G\) 是有限维连通实或复李群。以下条件等价：

\begin{enumerate}
\def\labelenumi{(\roman{enumi})}
\item
  L(G) 是半单的；
\item
  G 的根基是 \{e\}；
\item
  \(\mathbf{G}\) 的每个正规交换积分子群都等于 \(\{e\}\)。
\end{enumerate}

条件 (ii) 意味着 L(G) 的根基是 \(\{0\}\)，所以 (i) \(\Leftrightarrow\) (ii)（第一章第 6 节定理 1）。(i) 与 (iii) 的等价性由第 6 节第 6 号命题 14 得出。

定义 2. 若一个连通实或复李群是有限维的且满足命题 26 的条件，则称其为半单李群。

注 1. 设 G 是有限维连通实或复李群。若 G 不是半单的，则 G 有一个交换连通李子群 \(G'\)，它在每个连续自同构下都不变，且 \(G' \neq \{\varepsilon\}\)。事实上，令 n 为 L(G) 的最大幂零理想；则 \(n \neq \{0\}\)，相应的解析子群 N 是一个李子群，并在 G 的每个连续自同构下都不变（第 7 号命题 23）；N 的中心 G' 具有所需性质。

命题 27. 设 \(G\) 是有限维连通实或复李群。以下条件等价：

\begin{enumerate}
\def\labelenumi{(\roman{enumi})}
\item
  L(G) 是单的；
\item
  \(\mathbf{G}\) 的正规积分子群只有 \(\{e\}\) 和 \(\mathbf{G}\)，并且 \(\mathbf{G}\) 不是交换的。
\end{enumerate}

这由第 6 节第 6 号命题 14 得出。

定义 3. 若一个连通实或复李群是有限维的且满足命题 27 的条件，则称其为几乎单李群。

命题 28. 设 \(\mathbf{G}\) 是单连通实或复李群。以下条件等价：

\begin{enumerate}
\def\labelenumi{(\roman{enumi})}
\item
  G 是半单的；
\item
  G 同构于有限个几乎单群的乘积。
\end{enumerate}

若 G 是有限个几乎单李群的乘积，则 \(\mathrm{L}(\mathrm{G})\) 是有限个单李代数的乘积，因而是半单的。若 G 是半单的，则 \(\mathrm{L}(\mathrm{G})\) 同构于单李代数 \(L_{1},\ldots,L_{n}\) 的乘积。令 \(G_{i}\) 是李代数为 \(L_{i}\) 的单连通李群，因而它是几乎单的。于是 G 与 \(G_{1}\times\cdots\times G_{n}\) 都是单连通的，且具有同构的李代数，所以二者同构。

引理 1. 设 \(\mathbf{G}\) 是连通拓扑群，\(\mathbf{Z}\) 是其中心，\(\mathbf{Z}'\) 是 \(\mathbf{Z}\) 的离散子群。则 \(\mathbf{G}/\mathbf{Z}'\) 的中心是 \(\mathbf{Z}/\mathbf{Z}'\)。

设 \(y\) 是 G 的一个元素，其模 \(\mathbf{Z}'\) 的类是

G/Z' 的中心元素。令 \(\phi\) 为从 G 到 G 的映射 \(g \mapsto gyg^{-1}y^{-1}\)。则 \(\phi(G)\) 是连通的且包含于 Z' 中，所以 \(\phi(G) = \phi(\{e\}) = \{e\}\)。因此 \(y \in Z\)。

命题 29. 设 G 是半单连通实或复李群。

\begin{enumerate}
\def\labelenumi{(\roman{enumi})}
\item
  G = (G, G)。
\item
  G 的中心 Z 是离散的。
\item
  G/Z 的中心是 \{e\}。
\end{enumerate}

断言 (i) 由命题 4 的第 2 号推论和第一章第 6 节定理 1 得出。

断言 (ii) 由第 6 节第 4 号命题 10 的推论 4 以及第一章第 6 节第 1 号注 2 得出。

断言 (iii) 由 (ii) 和引理 1 得出。

命题 30. (i) 设 g 是半单实或复李代数。则 Int g 是 Aut g 的单位元分支。

\begin{enumerate}
\def\labelenumi{(\roman{enumi})}
\setcounter{enumi}{1}
\tightlist
\item
  设 \(\mathbf{G}\) 是半单连通实或复李群。\(\mathbf{G}\) 的伴随群是 \(\operatorname{Aut} \mathrm{L}(\mathrm{G})\) 的单位元分支。其中心退化为单位元。
\end{enumerate}

g 的每个导子都是内导子（第一章第 6 节命题 1 的推论 3），所以 L(Int g) = L(Aut g)，这就证明了 (i)。(ii) 的第一个断言由 (i) 得出，第二个断言由命题 29(iii) 以及第 6 节第 4 号命题 10 的推论 4(ii) 得出。

注 2. 设 g 是复半单李代数，\(g_{0}\) 是其底层实李代数。则 \(\operatorname{Aut}(g)\) 在 \(\operatorname{Aut}(g_{0})\) 中是开的，因为 \(\operatorname{Int}(g_{0}) \subset \operatorname{Aut}(g)\)。

命题 31. 设 G 是有限维单连通实或复李群，R 是其根基。存在 G 的一个半单单连通李子群 S，使得作为李群，G 是 S 与 R 的半直积。若 S' 是 G 的一个半单积分子群，则存在 L(G) 的幂零根基中的 x，使得

\[
(\operatorname{Ad} \exp x) (\mathrm{S} ^ {\prime}) \subset \mathrm{S}.
\]

这由第 6 节第 6 号命题 14 的推论 1 以及第一章第 6 节定理 5 和推论 1 得出。

引理 2. 设 G 是一个群（分别地，拓扑群），G' 是 G 的正规子群，V 是交换域 k（分别地，K）上的有限维向量空间，\(\rho\) 是 G 在 V 上的线性表示（分别地，连续线性表示），且 \(\rho' = \rho|G'\)。

\begin{enumerate}
\def\labelenumi{(\roman{enumi})}
\item
  若 \(\rho\) 是半单的，则 \(\rho'\) 是半单的。
\item
  若 \(\rho'\) 是半单的，且每个有限维线性 \(k\)-表示（分别地，连续线性 K-表示）都是 \(\mathrm{G} / \mathrm{G}'\)（分别地，\(\mathrm{G} / \overline{\mathrm{G}}'\)）的半单表示，则 \(\rho\) 是半单的。
\end{enumerate}

设 \(\rho\) 是半单的；我们证明 \(\rho'\) 是半单的。只需考虑 \(\rho\) 是单的情形。令 \(V'\) 为一个极小非零 \(G'\)-子模，它是 \(V\) 的子模。对所有 \(g \in G\)，\(\rho(G') \rho(g)V' = \rho(g)\rho(G')V' = \rho(g)V'\)，换言之，\(\rho(g)\mathrm{V}'\) 在 \(\rho(\mathrm{G}')\) 下稳定；若 \(\mathrm{V}''\) 是 \(\mathrm{G}'\) 的一个子模且包含于 \(\rho(g)\mathrm{V}'\)，则 \(\rho(g)^{-1}\mathrm{V}''\) 是 \(\mathrm{G}'\) 的一个子模且包含于 \(\mathrm{V}'\)，因而 \(\mathrm{V}''\) 等于 \(\{0\}\) 或 \(\rho(g)\mathrm{V}'\)。所以，对所有 \(g \in \mathrm{G}\)，\(\rho(g)\mathrm{V}'\) 都是一个单的 \(\mathrm{G}'\)-模。但 \(\sum_{g \in \mathrm{G}} \rho(g)\mathrm{V}'\) 是 \(\mathrm{V}\) 的非零子-G-模，从而 \(\mathrm{V} = \sum_{g \in \mathrm{G}} \rho(g)\mathrm{V}'\)。因此 \(\rho'\) 是半单的。

设 \(\rho'\) 是半单的。令 \(W\) 是 \(V\) 的一个非零子-G-模。由于 \(\rho'\) 是半单的，存在一个 \(f_0\)，它是从 \(V\) 到 \(W\) 的投影算子，并且与 \(\rho'(G)\) 对易。令 \(E\) 为满足以下条件的集合：其中的 \(f \in \mathcal{L}(V, V)\) 与 \(\rho(G')\) 对易，将 \(V\) 映入 \(W\)，且其在 \(W\) 上的限制是伸缩变换；对 \(f \in E\)，令 \(\alpha(f)\) 表示伸缩变换 \(f|W\) 的比值。于是 \(f_0 \in E\)，且 \(\alpha(f_0) = 1\)。显然，\(\alpha\) 是 \(E\) 上的线性形式。令 \(F = \text{Ker } \alpha\)，它是 \(E\) 的一个超平面。对 \(f \in E\) 和 \(g \in G\)，记 \(\sigma(g)f = \rho(g) \circ f \circ \rho(g)^{-1}\)；则 \(\sigma(g)f\) 将 \(V\) 映入 \(W\)，且其在 \(W\) 上的限制是比值为 \(\alpha(f)\) 的伸缩变换；若 \(g' \in G'\)，则

\[
\begin{array}{r l} \sigma (g) f \circ \rho (g ^ {\prime}) & = \rho (g) \circ f \circ \rho (g) ^ {- 1} \circ \rho (g ^ {\prime}) \\ & = \rho (g) \circ f \circ \rho (g ^ {- 1} g ^ {\prime} g) \circ \rho (g ^ {- 1}) \\ & = \rho (g) \circ \rho (g ^ {- 1} g ^ {\prime} g) \circ f \circ \rho (g ^ {- 1}) \\ & = \rho (g ^ {\prime}) \circ \rho (g) \circ f \circ \rho (g ^ {- 1}) \\ & = \rho (g ^ {\prime}) \circ \sigma (g) f. \end{array}
\]

因此 \(\sigma(g)f\in\mathbf{E}\)。于是 \(\sigma\) 是 G 在 E 上的线性 \(k\)-表示（分别地，连续线性 K-表示），并使 F 稳定。于是 \(\sigma(g)=\mathrm{Id}_{\mathbf{E}}\) 当 \(g\in\mathbf{G}'\) 时成立，在拓扑情形下进而当 \(g\in\overline{\mathbf{G}'}\) 时也成立。假设每个有限维 \(k\)-线性表示（分别地，连续 K-线性表示）作用在 G/G'（分别地，G/ \(\overline{\mathbf{G}'}\)）上都是半单的。于是 E 中存在 F 的一个在 G 下稳定的补空间。换言之，存在 \(f\in\mathbf{E}\)，使得 \(\alpha(f)=1\)，且 \(f\) 在 G 下不变。于是它是从 V 到 W 的投影算子，并且对 \(g\in\mathbf{G}\)，\(\rho(g)\circ f\circ\rho(g^{-1})=f\)，即 \(f\) 与 \(\rho(\mathbf{G})\) 对易。因此 \(\rho\) 是半单的。

定理 1. 设 G 是有限维实或复李群，\(G_{0}\) 是其单位元分支，R 是其根基，r 是 L(G) 的根基；假设 \(G/G_{0}\) 是有限的。设 \(\rho\) 是 G 的有限维解析线性表示。以下条件等价：

\begin{enumerate}
\def\labelenumi{(\roman{enumi})}
\item
  \(\rho\) 是半单的；
\item
  \(\rho|G_{0}\) 是半单的；
\item
  \(\rho|R\) 是半单的；
\item
  L(ρ) 是半单的；
\item
  L(ρ)\textbar r 是半单的。
\end{enumerate}

(i) \(\Leftrightarrow\) (ii) 由引理 2 和《积分》第七章第 3 节命题 1 得出。(ii) \(\Leftrightarrow\) (iv) 以及 (iii) \(\Leftrightarrow\) (v) 由第 6 节第 5 号命题 13 的推论 2 得出。(iv) \(\Leftrightarrow\) (v) 由第一章第 6 节定理 4 得出。

推论 1. 设 \(\rho, \rho_{1}, \rho_{2}\) 是 G 的有限维半单解析线性表示，n 是一个整数 \(\geqslant 0\)。则 \(\rho_{1} \otimes \rho_{2}\)、\(T^{n}\rho\)、\(S^{n}\rho\)、\(\bigwedge^{n}\rho\)（附录）都是半单的。

\(\rho_{1} \otimes \rho_{2}\) 的半单性由定理 1 和第一章第 6 节定理 4 的推论 1 得出。\(T^{n}_{\rho}, S^{n}_{\rho}, \bigwedge^{n}_{\rho}\) 的半单性由 \(\rho_{1} \otimes \rho_{2}\) 的半单性得出。

我们将在稍后看到，若 \(k\) 是特征为 0 的交换域，\(\Gamma\) 是一个群，且 \(\rho_1\) 和 \(\rho_2\) 是有限维半单线性 \(k\)-表示，其群为 \(\Gamma\)，则 \(\rho_1 \otimes \rho_2\) 是半单的。

推论 2. 设 \(\rho\) 是 \(G\) 在向量空间 \(V\) 上的有限维半单解析线性表示，\(S\) 是 \(V\) 的对称代数，\(S^{\mathfrak{G}}\) 是 \(S\) 的由在 \((\mathsf{S}(\rho))(G)\) 下不变的元素组成的子代数。则 \(S^{\mathfrak{G}}\) 是有限生成代数。

这由定理 1、第一章第 6 节定理 6(a) 以及《交换代数》第五章第 1 节定理 2 得出。

推论 3. 设 G 是实或复李群，\(G_{0}\) 是其单位元分支。假设 \(G_{0}\) 是半单的且 \(G/G_{0}\) 是有限的。则 G 的每个有限维解析线性表示都是半单的。

命题 32. 设 G 是有限维连通实李群。假设 L(G) 是约化的。以下条件等价：

\begin{enumerate}
\def\labelenumi{(\roman{enumi})}
\item
  \(\mathbf{G} / \overline{\mathbf{D}^1}\mathbf{G}\) 是紧的；
\item
  （分别地，(ii')）\(G\) 在复（分别地，实）向量空间上的每个有限维解析线性表示都是半单的。
\item
  \(\Rightarrow\) (ii')：假设 \(G / \overline{D}^1 G\) 是紧的。则 \(G / \overline{D}^1 G\) 在有限维实向量空间上的每个连续线性表示都是半单的（《积分》第七章第 3 节命题 1）。设 \(\rho\) 是 \(G\) 在实向量空间上的有限维解析线性表示。则 \(\rho |D^1 G\) 是解析的，\(D^1 G\) 是半单的（第一章第 6 节命题 5），所以 \(\rho |D^1 G\) 是半单的（定理 1 的推论 3）。因此 \(\rho\) 是半单的（引理 2）。
\end{enumerate}

同理可见 (i) \(\Rightarrow\) (ii)。

\begin{enumerate}
\def\labelenumi{(\roman{enumi})}
\tightlist
\item
  \(\Rightarrow\) (ii)：假设 \(G / \overline{D^1} G\) 不是紧的，因而同构于形如 \(\mathbf{R}^p\times \mathbf{T}^q\) 的群，其中 \(p > 0\)（第 6 节第 4 号命题 11(ii)）。于是存在从 \(G / \overline{D^1} G\) 到 \(\mathbf{R}\) 的满射态射，从而存在从 \(G\) 到 \(\mathbf{R}\) 的满射态射。映射
\end{enumerate}

\[
g \mapsto \sigma (g) = \left( \begin{array}{c c} 1 & 0 \\ \rho (g) & 1 \end{array} \right)
\]

是 G 在 \(R^{2}\) 上的解析线性表示，但不是半单的，因为在 \(\mathbf{R}^2\) 的一维向量子空间中，唯一在 \(\sigma(G)\) 下稳定的是 \(\mathbf{R}(0,1)\)。

同理可见 (ii) \(\Rightarrow\) (i)。

命题 33. 设 \(G\) 是连通分支数有限的有限维复李群，\(\rho\) 是 \(G\) 的有限维解析线性表示，\(G'\) 是实李群 \(G\) 的一个积分子群，且 \(L(G')\) 在 \(L(G)\) 上作为 \(C\)-向量空间生成。则 \(\rho\) 为半单的充要条件是 \(\rho |G'\) 为半单的。

令 \(\rho' = \rho|G'\)。\(\rho\)（分别地，\(\rho'\)）为半单的充要条件是 \(L(\rho)\)（分别地，\(L(\rho')\)）为半单的（定理 1）。令 \(V\) 为 \(\rho\) 的空间。\(V\) 的一个向量子空间在 \(L(\rho)(L(G))\) 下稳定的充要条件是它在 \(L(\rho')(L(G'))\) 下稳定。因此命题得证。

